% \def\hpitr{1}
\ifx\hpitr\undefined
\documentclass{article}
\setlength{\topmargin}{-0.25 in}
\setlength{\textwidth}{6.5in}
\setlength{\textheight}{8.5 in}
\setlength{\oddsidemargin}{0in}
\setlength{\evensidemargin}{0in}
\setlength{\parindent}{0.0in}
\setlength{\parskip}{7pt plus 1pt}
\else
\documentclass[english,trtype=singlearticle]{hpitr}
\fi

\usepackage{textcomp}
\usepackage{bcprules}
\usepackage{xcolor}				% Allow for color (annotations)
\usepackage{stmaryrd}
\usepackage{amssymb}
\usepackage{hyperref}
\usepackage{xspace}

\def\sectionautorefname{Section}
\def\subsectionautorefname{Section}
\def\subsubsectionautorefname{Section}
\def\tableautorefname{Table}
\def\figureautorefname{Figure}

\def\tdm#1{\textcolor{blue}{\textsf{TDM: #1}}}
\def\ab#1{\textcolor{red}{\textsf{AB: #1}}}
\def\tf#1{\textcolor{purple}{\textsf{TF: #1}}}

\newtheorem{Theorem}{Theorem}
\newtheorem{Corollary}{Corollary}
\newtheorem{Lemma}{Lemma}
\newtheorem{Claim}{Claim}
\newtheorem{Proposition}{Proposition}
\newtheorem{Definition}{Definition}

\newenvironment{proof}{\noindent\textbf{Proof.~~}}
{\hfill\ensuremath{\square}\\}

\newcommand{\maybe}{\textbf{\footnotesize{?}}\kern-1pt}

\newcommand{\Euid}{\texttt{\upshape E}}
\newcommand{\Huid}{\texttt{\upshape H}}
\newcommand{\Deref}{\texttt{\upshape $\eta$}}
\newcommand{\Cuid}{\texttt{\upshape C}}
\newcommand{\Iuid}{\texttt{\upshape I}}

\newcommand{\eval}[3]{\texttt{#1} $\vdash$ \texttt{#2} $\Downarrow$ \texttt{#3}}
\newcommand{\gstay}[3]{stay(\texttt{#1}, \texttt{#2}) $=$ \texttt{#3}}
\newcommand{\stayPrefix}[4]{stayPrefix(\texttt{#1}, \texttt{#2}, \texttt{#3}) $=$ \texttt{#4}}
\newcommand{\phase}[5]{solve(\texttt{#1}, \texttt{#2}, \texttt{#3}, \texttt{#4}) $=$ \texttt{#5}}
\newcommand{\stay}[2]{stay(\texttt{#1}) $=$ \texttt{#2}}
\newcommand{\istay}[2]{i-stay(\texttt{#1}) $=$ \texttt{#2}}
\newcommand{\toSolver}[2]{\texttt{E} $\vdash$ \texttt{#1} $\leadsto$ $#2$}
\newcommand{\fromSolver}[2]{$#1$ $\leadsto$ \texttt{#2}}
\newcommand{\csolve}[2]{$#1$ $\models$ $#2$}
\newcommand{\solve}[2]{\texttt{#1} $\models$ \texttt{#2}}
\newcommand{\cfail}[1]{$\not\models$ $#1$}
\newcommand{\fail}[1]{$\not\models$ \texttt{#1}}
\newcommand{\step}[5]{\texttt{<#1|#2|#3>} $\longrightarrow$
  \texttt{<#4|#5>}}
\newcommand{\tc}[3]{\texttt{#1} $\vdash$ \texttt{#2} : \texttt{#3}}
\newcommand{\tcC}[2]{\texttt{#1} $\vdash$ \texttt{#2}}

\newcommand{\evalH}[4]{\texttt{#1;#2} $\vdash$ \texttt{#3} $\Downarrow$ \texttt{#4}}
\newcommand{\solveH}[3]{\texttt{#1;#2} $\models$ \texttt{#3}}
% \newcommand{\solveH}[3]{\texttt{#1;#2} $\models_{E,H}$ \texttt{#3}}
% \newcommand{\solveH}[3]{\texttt{#1} $\models_{E,H}$ \texttt{#3}}
\newcommand{\failH}[1]{$\not\models_{E,H}$ \texttt{#1}}
\newcommand{\stepH}[9]{\texttt{<#1|#2|#3|#4|#5>} $\longrightarrow$
  \texttt{<#6|#7|#8|#9>}}
\newcommand{\tcH}[4]{\texttt{#1;#2} $\vdash$ \texttt{#3} {\upshape :} \texttt{#4}}
\newcommand{\tcHC}[3]{\texttt{#1;#2} $\vdash$ \texttt{#3}}
\newcommand{\toSolverH}[2]{\texttt{E,H} $\vdash$ \texttt{#1} $\leadsto$ $#2$}
\newcommand{\fromSolverH}[3]{$#1$ $\leadsto$ \texttt{#2,#3}}

\newcommand{\toSolverOO}[6]{\texttt{#1,#2} $\vdash$ \texttt{#3} $\leadsto$ \texttt{#4,#5} $\mid$ $#6$}

\newcommand{\inferrule}[3]{\infrule[#1]{\mbox{#2}}{\mbox{#3}}}
\newcommand{\inferax}[2]{\infax[#1]{\mbox{#2}}}

\newcommand{\judgment}[1]{\begin{flushleft} \fbox{#1}\end{flushleft}}

\newcommand{\receq}{\stackrel{\texttt{r}}{\texttt{=}}}

\newcommand{\lookup}[3]{\textit{lookup}(\texttt{#1},\texttt{#2}) = (\texttt{x$_1$} $\cdots$ \texttt{x$_n$},\texttt{#3})}


\newcommand{\combine}[4]{\texttt{#1;#2} $\vdash$ \texttt{#3} $\leadsto$ \texttt{#4}}
\newcommand{\inline}[9]{\texttt{<#1,#2,#3,#4,#5,#6>} $\leadsto$ \texttt{<#7,#8,#9>}}


\newcommand\enter[9]{%
\def\argA{#1}\def\argB{#2}\def\argC{#3}\def\argD{#4}\def\argE{#5}%
\def\argF{#6}\def\argG{#7}\def\argH{#8}\def\argI{#9}%
\entercontinued
}
\newcommand\entercontinued[2]{%
  \textit{enter}(\texttt{\argA},\texttt{\argB},\texttt{\argC},\texttt{\argD},\texttt{
    \argE},\texttt{\argF},\texttt{x$_1$} $\cdots$ \texttt{x$_n$},\texttt{e$_1$} $\cdots$
  \texttt{e$_n$}) = (\texttt{\argG},\texttt{\argH},\texttt{\argI},\texttt{#1},\texttt{#2})
}

\newcommand\stepHO[9]{%
\def\argA{#1}\def\argB{#2}\def\argC{#3}\def\argD{#4}\def\argE{#5}%
\def\argF{#6}\def\argG{#7}\def\argH{#8}\def\argI{#9}%
\stepHOcontinued
}
\newcommand{\stepHOcontinued}[2]{%
  \texttt{<\argA|\argB|\argC|\argD|\argE|\argF>} $\longrightarrow$ \texttt{<\argG|\argH|\argI|#1|#2>}}


\newcommand\evalBig[9]{%
\def\argA{#1}\def\argB{#2}\def\argC{#3}\def\argD{#4}\def\argE{#5}%
\def\argF{#6}\def\argG{#7}\def\argH{#8}\def\argI{#9}%
\evalBigcontinued
}
\newcommand{\evalBigcontinued}[2]{%
  \texttt{<\argA|\argB|\argC|\argD|\argE|\argF>} $\Downarrow$ \texttt{<\argG|\argH|\argI|#1|#2>}}

\newcommand{\resolve}[4]{\texttt{#3;#4} $\models_\texttt{#1}$ \texttt{#2}}

\newcommand{\E}{$\mathbb{E}$}
\newcommand{\Ccal}{$\mathbb{C}$}
\newcommand{\Ical}{$\mathbb{I}$}
\renewcommand{\H}{$\mathbb{H}$}
\newcommand{\reinline}[6]{\texttt{<#1,#2,#3,#4>} $\leadsto$ \texttt{<#5,#6>}}

\ifx\C\undefined
\newcommand{\C}{\Ccal}
\else
\renewcommand{\C}{\Ccal}
\fi
\newcommand{\I}{\Ical}

\usepackage{listings}
\usepackage{needspace}
\newcounter{examples}
\lstnewenvironment{example}[1][]%
{\stepcounter{examples}%
  \needspace{2\baselineskip}
  \vspace{2\topsep}%
  \marginpar{\raggedleft\textit{Test \arabic{examples}}}%
  \vspace{-2\topsep}%
  \lstset{%
    basicstyle=\ttfamily,%
    frame=l,%
    #1}}{}

% \newcommand{\example}[1]{%
% % \leavevmode %
% \vspace{2\topsep}%
% \marginpar{\raggedleft\textit{Test #1}}%
% \vspace{-2\topsep}%
% % \vspace{-3pt}%
% }

\usepackage{etoolbox}
\makeatletter
\preto{\@verbatim}{\topsep=0pt \partopsep=0pt}
\makeatother

\sloppy

\begin{document}
\reversemarginpar

\title{Checks and Balances -- Constraint Solving without Surprises \\
in Object-Constraint Programming Languages: \\
Full Formal Development}
\author{Tim Felgentreff, Todd Millstein, Alan Borning, and Robert Hirschfeld}

\date{VPRI Technical Report TR-2015-001, July 2015}


\maketitle

\section{Two Key Theorems}

This technical report is intended to accompany the conference paper 
``Constraint Solving without Surprises in Object-Constraint Programming
Languages,'' and presents the proofs for two theorems that capture key
properties of the Babelsberg/UID core language.  These theorems are
stated (without the proofs) in Section 3.3 of that paper.

%% used to customize a sentence about where assumptions about
%% constraint solver properties are described
\newcommand{\where}{as described in the main conference paper}

%% This is the proofs of the two main theorems for Babelsberg/UID,
%% broken out into a separate file for inclusion in both the TR to
%% accompany the OOPSLA paper, and also the step-by-step development
%% tech report.

We
assume that a given configuration {\tt \Euid|\Huid|C|I} is {\em well
  formed}, meaning that it satisfies these sanity conditions:
\begin{itemize}
  \item any object reference appearing in the range of \Euid\ also
    appears in the domain of \Huid  
 \item any object reference appearing in the range of \Huid\ also
    appears in the domain of \Huid
 \item for all variables {\tt x} in the domain of \Euid\ there is some
   type {\tt T} such that \tcH{\Euid}{\Huid}{x}{T}
 \item \tcHC{\Euid}{\Huid}{C}
 \item \evalH{\Euid}{\Huid}{I}{true}
\end{itemize}
Well-formedness follows from the assumptions made on the constraint
solver \where.
% The first three assumptions above ensure some natural sanity conditions.  
% The last condition above is an artifact of our formalism's simple type
% system, which lacks recursive types.  Therefore it is easy to show
% that recursive objects can never be created, and one lemma below leverages
% this knowledge in order to argue that every value has a type.

The
first theorem formalizes the idea that any solution to a value
constraint preserves the structures of the objects on the
environment and heap:

\begin{Theorem} (Structure Preservation)
 If \emph{\stepH{\Euid}{\Huid}{\C}{\I}{s}{\Euid$'$}{\Huid$'$}{\C$'$}{\I$'$}}
 and {\tt s} either has the form {\tt once C$_0$} or {\tt always
   C$_0$} and \emph{\tcHC{\Euid}{\Huid}{C$_0$}} and
\emph{\tcH{\Euid}{\Huid}{L}{T}}, then
  \emph{\tcH{\Euid$'$}{\Huid$'$}{L}{T}}. 
\end{Theorem}
\begin{proof}
If {\tt s} has the form {\tt once C$_0$} then the result follows by
Lemma~\ref{lem:key}.  If {\tt s} has the form {\tt always C$_0$}, then
since
\stepH{\Euid}{\Huid}{\C}{\I}{s}{\Euid$'$}{\Huid$'$}{\C$'$}{\I$'$} and {\tt C$_0$} is not an identity test, by
rule {\sc S-Always} we have 
\stepH{\Euid}{\Huid}{\C}{\I}{once C$_0$}{\Euid$'$}{\Huid$'$}{\C}{\I}.
Then again the result follows from Lemma~\ref{lem:key}.
\end{proof}

\begin{Lemma}
\label{lem:key}
 If \emph{\stepH{\Euid}{\Huid}{\C}{\I}{once
   C$_0$}{\Euid$'$}{\Huid$'$}{\C$'$}{\I$'$}}
and \emph{\tcHC{\Euid}{\Huid}{C$_0$}} and
\emph{\tcH{\Euid}{\Huid}{L}{T}}, then
  \emph{\tcH{\Euid$'$}{\Huid$'$}{L}{T}}. 
\end{Lemma}
\begin{proof}
Since \stepH{\Euid}{\Huid}{\C}{\I}{once
   C$_0$}{\Euid$'$}{\Huid$'$}{\C$'$}{\I$'$} and  {\tt C$_0$} is not an
 identity test, by {\sc S-Once} we have that 
\tcHC{\Euid}{\Huid}{C$_0$} and
    \phase{\Euid}{\Huid}{\C\,$\wedge$\,C$_0$}{required}{\Euid$'$;\Huid$'$}.
By the assumption of well-formedness we have \tcHC{\Euid}{\Huid}{C} so
by {\sc T-Conjunction} also \tcHC{\Euid}{\Huid}{C\,$\wedge$\,C$_0$}. 
Therefore the result follows from Lemma~\ref{lem:keyAux}.
\end{proof}

\begin{Lemma}
\label{lem:keyAux}
 If \emph{\tcHC{\Euid}{\Huid}{C}} and
 \emph{\phase{\Euid}{\Huid}{\C}{required}{\Euid$'$;\Huid$'$}}  and
\emph{\tcH{\Euid}{\Huid}{L}{T}}, then
  \emph{\tcH{\Euid$'$}{\Huid$'$}{L}{T}}. 
\end{Lemma}
\begin{proof}
Since \phase{\Euid}{\Huid}{\C}{required}{\Euid$'$;\Huid$'$} by
{\sc Solve} we have that 
    \gstay{\Euid}{required}{C$_{E_s}$} and
    \gstay{\Huid}{required}{C$_{H_s}$} and
    \combine{\Euid}{\Huid}{\C}{C$'$} and
    \solveH{\Euid$'$}{\Huid$'$}{$($C$'$ $\wedge$ C$_{E_s}$ $\wedge$
      C$_{H_s}$$)$}.
We prove this lemma by structural induction on {\tt T}.  

By Lemma~\ref{lem:typeImpliesValue} there exists a value {\tt v}
such that \evalH{\Euid}{\Huid}{L}{v} and \tcH{\Euid}{\Huid}{v}{T}.
Then by Lemma~\ref{lem:aux} there exists a value {\tt
  v$'$} such that
  \evalH{\Euid$'$}{\Huid$'$}{L}{v$'$}.  Furthermore, if {\tt v}
  is an object reference {\tt r}, then also {\tt v$'$} = {\tt r}.
  Case analysis on the form of {\tt v$'$}:
 \begin{itemize}
\item Case {\tt v$'$} is a constant {\tt c$'$}.  Then also {\tt v} is
  a constant {\tt c}.  Since \tcH{\Euid}{\Huid}{v}{T}, the last rule
  in this derivation must be {\sc T-Constant} so {\tt T} is {\tt
    PrimitiveType} and the result follows from {\sc T-Constant}.
  \item Case {\tt v$'$} is a reference {\tt r$'$}.  Since 
\evalH{\Euid$'$}{\Huid$'$}{L}{v$'$} by 
  Lemma~\ref{lem:valueImpliesType} there is a type {\tt T$'$} such
  that
  \tcH{\Euid$'$}{\Huid$'$}{L}{T$'$}   and
  \tcH{\Euid$'$}{\Huid$'$}{r$'$}{T$'$}.   
So it suffices to show that {\tt
    T} = {\tt T$'$}.  Case analysis on the form of
  {\tt v}:
\begin{itemize}
\item Case {\tt v} is a constant {\tt c}.  Since \tcH{\Euid}{\Huid}{v}{T}, the
    last rule in this derivation must be {\sc T-Constant}, so 
   we have {\tt T} = {\tt PrimitiveType}.  Case analysis on the last
   rule in the derivation of \evalH{\Euid}{\Huid}{L}{v}:
  \begin{itemize}
  \item Case {\sc E-Variable}.  Then {\tt L} is a variable {\tt x} and
  \Euid({\tt x}) = {\tt c}.  Since \gstay{\Euid}{required}{C$_{E_s}$}, by
  {\sc StayEnv} we have that {\tt C$_{E_s}$} includes a conjunct {\tt
    C$_x$} such that \gstay{x=c}{required}{C$_x$}, and by {\sc
    StayConst}
{\tt C$_x$} has the form {\tt weak x=c}.  

Since
\evalH{\Euid$'$}{\Huid$'$}{L}{r$'$}, the last rule in this derivation
must be {\sc E-Variable}, so we have \Euid$'$({\tt x}) = {\tt r$'$}.  
We know that assigning {\tt x} to {\tt c} satisfies the weak stay
constraint above.
Therefore there must be some constraint on {\tt x} in {\tt C}
that causes the change in value for {\tt x} from {\tt c} to {\tt
  r$'$}.  
  \item Case {\sc E-Field}.  Then {\tt L} has the form  {\tt L$_0$.l$_i$}
  and \evalH{\Euid}{\Huid}{L$_0$}{r$_0$} and
  \texttt{\Huid}(\texttt{r$_0$}) =
  \texttt{\{l$_1$:v$_1$,$\ldots$,l$_n$:v$_n$\}} and
  1 $\leq$ $i$ $\leq$ $n$ and {\tt r$'$} = {\tt v$_i$}. 
 Since \gstay{\Huid}{required}{C$_{H_s}$}, by
  {\sc StayHeap} we have {\tt C$_{H_s}$} includes a conjunct {\tt
    C$_{r_0}$} such that \gstay{r$_0$=\{l$_1$:v$_1$,$\ldots$,l$_n$:v$_n$\}}{required}{C$_{r_0}$}.
By {\sc StayObject} {\tt C$_{r_0}$} has the form {\tt (required \Deref(r$_0$)=\{l$_1$:x$_1$,$\ldots$,l$_n$:x$_n$\}) $\wedge$
    C$_{1}$ $\wedge$
    $\cdots$ $\wedge$
    C$_{n}$}, where {\tt x$_1$} $\ldots$ {\tt x$_n$} are fresh
  variables and \gstay{x$_i$=v$_i$}{required}{C$_i$}.  Then by {\sc
    StayConst}
{\tt C$_i$} has the form {\tt weak x$_i$=c}.  

Since \evalH{\Euid}{\Huid}{L$_0$}{r$_0$}, by
  Lemma~\ref{lem:aux} we have \evalH{\Euid$'$}{\Huid$'$}{L$_0$}{r$_0$}.
Since
\evalH{\Euid$'$}{\Huid$'$}{L}{r$'$}, the last rule in this derivation
must be {\sc E-Field}, so we have that \Huid$'$({\tt r$_0$}) =
\texttt{\{l$_1$:v$_1'$,$\ldots$,l$_n$:v$_n'$\}}  where {\tt v$_i'$}, and
hence {\tt x$_i$} is
{\tt r$'$}.
We know that assigning {\tt x$_i$} to {\tt c} satisfies the weak stay
constraint above.
Therefore there must be some constraint on an l-value of the form {\tt
  L$_{00}$.l$_i$} in {\tt C}, where \evalH{\Euid}{\Huid}{L$_{00}$}{r$_0$},
that causes the change in value for {\tt x$_i$} from {\tt c} to {\tt
  r$'$}.  
\end{itemize}

Therefore in either case, there must be some constraint on an l-value
{\tt L$''$} in {\tt C}, where
\evalH{\Euid}{\Huid}{L$''$}{c} but
\evalH{\Euid$'$}{\Huid$'$}{L$''$}{r$'$}.
By Lemma~\ref{lem:valueImpliesType} there is some {\tt T$''$} such
that \tcH{\Euid}{\Huid}{L$''$}{T$''$}
and \tcH{\Euid}{\Huid}{c}{T$''$}.  The last rule in the derivation of
\tcH{\Euid}{\Huid}{c}{T$''$} must be {\sc T-Constant} so {\tt T$''$}
is {\tt PrimitiveType}.
We argue a contradiction by case analysis of the immediate parent
expression of any occurrence of {\tt L$''$} in {\tt C}.  We are given \tcHC{\Euid}{\Huid}{C} so this parent
expression must also be well typed.
  \begin{itemize}
   \item Case {\tt L$''$.l}.  Then by {\sc T-Field} {\tt L$''$} must have a
     record type, contradicting the fact that \tcH{\Euid}{\Huid}{L$''$}{PrimitiveType}.
  \item Case {\tt L$''$} is an immediate subexpression of a $\oplus$,
    $\olessthan$, or $\ovee$ operation.  But these operations only
    apply to primitives, so the solver cannot satisfy them by
    assigning {\tt x} to {\tt r$'$}.
  \item Case {\tt $\rho$ L$''$}.  This constraint is only satisfied if
    {\tt x} is a boolean, so the solver will not assign {\tt x} to
    {\tt r$'$}.
  \end{itemize}

  \end{itemize}

\item Case {\tt v} is a reference {\tt r}.  Then {\tt r} = {\tt
    r$'$}.  Since \tcH{\Euid}{\Huid}{v}{T}, by {\sc T-Ref} we have
  that \texttt{\Huid}(\texttt{r})=\texttt{\{l$_1$:v$_1$,$\ldots$,l$_n$:v$_n$\}}
    and
  \tcH{\Euid}{\Huid}{v$_1$}{T$_1$} $\cdots$
  \tcH{\Euid}{\Huid}{v$_n$}{T$_n$} and {\tt T} is
  \texttt{\{l$_1$:T$_1$,$\ldots$,l$_n$:T$_n$\}}.  By {\sc T-Field} we
  have \tcH{\Euid}{\Huid}{L.l$_i$}{T$_i$} for each $1 \leq i \leq n$,
  so by induction \tcH{\Euid$'$}{\Huid$'$}{L.l$_i$}{T$_i$} for each $1
  \leq i \leq n$.

 Since \gstay{\Huid}{required}{C$_{H_s}$}, by
  {\sc StayHeap} we have that {\tt C$_{H_s}$} includes a conjunct {\tt
    C$_r$} such that \gstay{r=\{l$_1$:v$_1$,$\ldots$,l$_n$:v$_n$\}}{required}{C$_r$}.
By {\sc StayObject} {\tt C$_r$} has the form {\tt (required \Deref(r)=\{l$_1$:x$_1$,$\ldots$,l$_n$:x$_n$\}) $\wedge$
    C$_{1}$ $\wedge$
    $\cdots$ $\wedge$
    C$_{n}$}, where {\tt x$_1$} $\ldots$ {\tt x$_n$} are fresh
  variables and \gstay{x$_i$=v$_i$}{required}{C$_i$}.  Therefore any solution to {\tt $($C$'$ $\wedge$ C$_{E_s}$ $\wedge$
      C$_{H_s}$$)$} must map {\tt r} to an object value of the form
    {\tt  \{l$_1$:v$_1'$,$\ldots$,l$_n$:v$_n'$\}} in {\tt \Huid$'$}.
    
Since \tcH{\Euid$'$}{\Huid$'$}{L.l$_i$}{T$_i$} for each $1
  \leq i \leq n$,  by Lemma~\ref{lem:typeImpliesValue} 
we have 
  \evalH{\Euid}{\Huid}{L.l$_i$}{v$_i''$}  and
  \tcH{\Euid}{\Huid}{v$_i''$}{T$_i$} for each $1 \leq i \leq n$.
  Since the last rule in each of these evaluation derivations must be
  {\sc E-Field} we have that {\tt v$_i''$} = {\tt v$_i'$} for each $1
  \leq i \leq n$.
  Finally, since \tcH{\Euid$'$}{\Huid$'$}{r}{T$'$},  the last rule in
  this derivation must be {\sc T-Ref}, so we
  have that {\tt T$'$} = {\tt  \{l$_1$:T$_1$,$\ldots$,l$_n$:T$_n$\}}.

 \end{itemize}
\end{proof}

\begin{Lemma}
\label{lem:typeImpliesValue}
If \emph{\tcH{\Euid}{\Huid}{L}{T}}, then there exists a value {\tt v}
such that \emph{\evalH{\Euid}{\Huid}{L}{v}} and \emph{\tcH{\Euid}{\Huid}{v}{T}}.
\end{Lemma}
\begin{proof}
By structural induction on {\tt L}:
\begin{itemize}
\item Case {\tt L} is a variable {\tt x}:  Then by {\sc T-Variable} we have
that \texttt{\Euid}(\texttt{x}) = \texttt{v} and
\tcH{\Euid}{\Huid}{v}{T}.  Finally by {\sc E-Var} we
have \evalH{\Euid}{\Huid}{x}{v}.

\item Case {\tt L} has the form {\tt L$'$.l$_i$}:  By {\sc T-Field} we
  have that \tcH{\Euid}{\Huid}{L$'$}{\{l$_1$:T$_1$,$\ldots$,l$_n$:T$_n$\}}
  and $1 \leq i \leq n$ and {\tt T} = {\tt T$_i$}.  By induction there
  exists a value {\tt v$'$} such that \evalH{\Euid}{\Huid}{L$'$}{v$'$}
  and \tcH{\Euid}{\Huid}{v$'$}{\{l$_1$:T$_1$,$\ldots$,l$_n$:T$_n$\}}.
  Case analysis of the derivation of
  \tcH{\Euid}{\Huid}{v$'$}{\{l$_1$:T$_1$,$\ldots$,l$_n$:T$_n$\}}:
  \begin{itemize}
    \item Case {\sc T-Constant}:  Then {\tt \{l$_1$:T$_1$,$\ldots$,l$_n$:T$_n$\}} = {\tt PrimitiveType}
      and we have a contradiction.
    \item Case {\sc T-Ref}:  Then we have {\tt v$'$} = {\tt r} and 
  \texttt{\Huid}(\texttt{r})=\texttt{\{l$_1$:v$_1$,$\ldots$,l$_n$:v$_n$\}}
    and
  \tcH{\Euid}{\Huid}{v$_i$}{T$_i$}.  Finally, by {\sc E-Field}
  we have \evalH{\Euid}{\Huid}{L$'$.l$_i$}{v$_i$}.
  %   \item Case {\tt T-Object}:  Then we have {\tt v$'$} = {\tt \{l$_1$:v$_1$,$\ldots$,l$_n$:v$_n$\}} and
  % \tcH{\Euid}{\Huid}{v$_i$}{T$_i$}.   Finally, by {\sc E-ValueField}
  % we have \evalH{\Euid}{\Huid}{L$'$.l$_i$}{v$_i$}.
   \end{itemize}
\end{itemize}
\end{proof}

\begin{Lemma}
\label{lem:aux}
If \emph{\tcHC{\Euid}{\Huid}{C}} and
 \emph{\phase{\Euid}{\Huid}{\C}{required}{\Euid$'$;\Huid$'$}}  and
\emph{\evalH{\Euid}{\Huid}{L}{v}}, then there exists a value {\tt
  v$'$} such that
  \emph{\evalH{\Euid$'$}{\Huid$'$}{L}{v$'$}}.  Furthermore, if {\tt v}
  is an object reference {\tt r}, then also {\tt v$'$} = {\tt r}.
\end{Lemma}
\begin{proof}
Since \phase{\Euid}{\Huid}{\C}{required}{\Euid$'$;\Huid$'$} by
{\sc Solve} we have that 
    \gstay{\Euid}{required}{C$_{E_s}$} and
    \gstay{\Huid}{required}{C$_{H_s}$} and
    \combine{\Euid}{\Huid}{\C}{C$'$} and
    \solveH{\Euid$'$}{\Huid$'$}{$($C$'$ $\wedge$ C$_{E_s}$ $\wedge$
      C$_{H_s}$$)$}.
 We proceed by structural induction on {\tt L}:
\begin{itemize}
\item Case {\tt L} is a variable {\tt x}.  Since
  \evalH{\Euid}{\Huid}{L}{v}, by {\sc E-Var} we have that \Euid({\tt
    x}) = {\tt v}.  Then since \gstay{\Euid}{required}{C$_{E_s}$}, by
  {\sc StayEnv} we have that {\tt C$_{E_s}$} includes a conjunct {\tt
    C$_x$} such that \gstay{x=v}{required}{C$_x$}.  Case analysis of
  the rule used in the derivation of \gstay{x=v}{required}{C$_x$}:
\begin{itemize}
 \item {\sc StayConst}:  Then 
{\tt C$_x$} has the form {\tt weak x=v} and {\tt v} is a constant {\tt
  c}.
  Therefore the variable {\tt x} appears in the constraint  {\tt $($C$'$ $\wedge$ C$_{E_s}$ $\wedge$
      C$_{H_s}$$)$}, so any solution to the constraint must include some value {\tt v$'$} for {\tt x} in
    {\tt \Euid$'$}, and the result follows by {\sc E-Var}.

 \item {\sc StayRef}:  Then 
{\tt C$_x$} has the form {\tt required x=v} and {\tt v} is an object reference {\tt
  r}.
  Therefore any solution to  {\tt $($C$'$ $\wedge$ C$_{E_s}$ $\wedge$
      C$_{H_s}$$)$} must map {\tt x} to {\tt r} in
    {\tt \Euid$'$}, and the result follows by {\sc E-Var}.
\end{itemize}
\item Case {\tt L} has the form {\tt L$'$.l$_i$}.  Then by {\sc
    E-Field} we have \evalH{\Euid}{\Huid}{L$'$}{r} and
  \texttt{\Huid}(\texttt{r}) =
  \texttt{\{l$_1$:v$_1$,$\ldots$,l$_n$:v$_n$\}} and
  1 $\leq$ $i$ $\leq$ $n$ and {\tt v} = {\tt v$_i$}.  By induction we have
\evalH{\Euid$'$}{\Huid$'$}{L$'$}{r}.  Since \gstay{\Huid}{required}{C$_{H_s}$}, by
  {\sc StayHeap} we have that {\tt C$_{H_s}$} includes a conjunct {\tt
    C$_r$} such that \gstay{r=\{l$_1$:v$_1$,$\ldots$,l$_n$:v$_n$\}}{required}{C$_r$}.
By {\sc StayObject} {\tt C$_r$} has the form {\tt (required \Deref(r)=\{l$_1$:x$_1$,$\ldots$,l$_n$:x$_n$\}) $\wedge$
    C$_{1}$ $\wedge$
    $\cdots$ $\wedge$
    C$_{n}$}, where {\tt x$_1$} $\ldots$ {\tt x$_n$} are fresh
  variables and \gstay{x$_i$=v$_i$}{required}{C$_i$}.  Therefore any solution to {\tt $($C$'$ $\wedge$ C$_{E_s}$ $\wedge$
      C$_{H_s}$$)$} must map {\tt r} to an object value of the form
    {\tt  \{l$_1$:v$_1'$,$\ldots$,l$_n$:v$_n'$\}} in {\tt \Huid$'$},
    where {\tt v$_i'$} is the value assigned to variable {\tt x$_i$}
    by the solution.
    Therefore by {\sc E-Field} we have
    \evalH{\Euid}{\Huid}{L$'$.l$_i$}{v$_i'$}.  Finally  suppose {\tt
      v$_i$} is an object reference {\tt r$_i$}.  Then by {\sc
      StayRef} {\tt C$_i$} has the form {\tt required x$_i$ = r$_i$} so any
    solution to the constraints must map  {\tt x$_i$} to {\tt r$_i$},
    so also {\tt v$_i'$} is {\tt r$_i$}.
\end{itemize}
\end{proof}



\begin{Lemma}
\label{lem:valueImpliesType}
If \emph{\evalH{\Euid}{\Huid}{L}{v}} then there is some type {\tt T}
such that \emph{\tcH{\Euid}{\Huid}{L}{T}} and \emph{\tcH{\Euid}{\Huid}{v}{T}}.
\end{Lemma}
\begin{proof}
By induction on the derivation of \evalH{\Euid}{\Huid}{L}{v}:
\begin{itemize}
\item Case {\sc E-Variable}.  Then {\tt L} is a variable {\tt x} and 
  \Euid({\tt x}) = {\tt v}.  Then by
  well formedness there is some type {\tt T}
such that \tcH{\Euid}{\Huid}{x}{T} and by {\tt T-Variable} also \tcH{\Euid}{\Huid}{v}{T}.
\item Case {\sc E-Field}.  Then {\tt L} has the form {\tt L$'$.l$_i$}
  and \evalH{\Euid}{\Huid}{L$'$}{r} and
  \texttt{\Huid}(\texttt{r}) =
  \texttt{\{l$_1$:v$_1$,$\ldots$,l$_n$:v$_n$\}} and
  1 $\leq$ $i$ $\leq$ $n$ and {\tt v} = {\tt v$_i$}.  By induction
  there is some type {\tt T$'$} such that
  \tcH{\Euid}{\Huid}{L$'$}{T$'$} and \tcH{\Euid}{\Huid}{r}{T$'$}.
  Then by {\sc T-Ref} {\tt T$'$} has the form 
  \texttt{\{l$_1$:T$_1$,$\ldots$,l$_n$:T$_n$\}} where
\tcH{\Euid}{\Huid}{v$_i$}{T$_i$}.  Then by {\sc T-Field} also 
\tcH{\Euid}{\Huid}{L}{T$_i$}.
\end{itemize}
\end{proof}



\begin{Lemma}
\label{lem:refWellFormed}
If \emph{\evalH{\Euid}{\Huid}{L}{r}} then {\tt r} is in the domain of \Huid.
\end{Lemma}
\begin{proof}
Case analysis on the last rule in the evaluation derivation.
\begin{itemize}
\item Case {\sc E-Variable}.  Then {\tt L} is a variable {\tt x} and
  \Euid({\tt x}) = {\tt r}.  Then the result follows from the
  assumption that \Euid\ and \Huid\ are well formed.
\item Case {\sc E-Field}.  Then {\tt L} has the form {\tt L$'$.l$_i$}
  and \evalH{\Euid}{\Huid}{L$'$}{r$'$} and
  \texttt{\Huid}(\texttt{r}) =
  \texttt{\{l$_1$:v$_1$,$\ldots$,l$_n$:v$_n$\}} and
  1 $\leq$ $i$ $\leq$ $n$ and {\tt r} = {\tt v$_i$}.  Then the result
  follows from the assumption that  \Euid\ and \Huid\ are well formed.
\end{itemize}
\end{proof}

% \begin{Lemma}
% \label{lem:cf}
% (Canonical Forms)
% \begin{itemize}
% \item If \emph{\tcH{\Euid}{\Huid}{v}{PrimitiveType}} then {\tt v} is a
% constant {\tt c}.
% \item If
%   \emph{\tcH{\Euid}{\Huid}{v}{\{l$_1$:T$_1$,$\ldots$,l$_n$:T$_n$\}}}
%   then 
% {\tt v} is an object reference {\tt r} such that \Deref({\tt
%        r}) has the form {\tt \{l$_1$:v$_1$,$\ldots$,l$_n$:v$_n$\}}.
% \end{itemize}
% \end{Lemma}
% \begin{proof}
% Immediate from the typing rules.
% \end{proof}

% \begin{Lemma}
% \label{lem:key}
%  If \stepH{\Euid}{\Huid}{\C}{\I}{once
%    C$_0$}{\Euid$'$}{\Huid$'$}{\C$'$}{\I$'$} and {\tt C$_0$} is not an
%  identity test and
% and \evalH{\Euid}{\Huid}{L}{v}, then:
% \begin{itemize}
% \item If \texttt{v} is a constant \texttt{c}, then
% \evalH{\Euid$'$}{\Huid$'$}{L}{c$'$} for some constant \texttt{c$'$}.  
% \item If \texttt{v} is an object \texttt{\{l$_1$:v$_1$,$\ldots$,l$_n$:v$_n$\}}, then
% \evalH{\Euid$'$}{\Huid$'$}{L}{\{l$_1$:v$_1'$,$\ldots$,l$_n$:v$_n'$\}} for some
% values \texttt{v$_1'$} $\cdots$ \texttt{v$_n'$}.
% \item If \texttt{v} is a reference \texttt{r} and \Deref(\texttt{r})=\texttt{\{l$_1$:v$_1$,$\ldots$,l$_n$:v$_n$\}}, then
% \evalH{\Euid$'$}{\Huid$'$}{L}{r} and \Huid$'$(\texttt{r})=\texttt{\{l$_1$:v$_1'$,$\ldots$,l$_n$:v$_n'$\}} for some
% values \texttt{v$_1'$} $\cdots$ \texttt{v$_n'$}.
% \end{itemize}
% \end{Lemma}
% \begin{proof}
% Since  \tcHC{\Euid}{\Huid}{C$_0$} we have that \texttt{C$_0$} is a value
% constraint rather than an identity constraint.  Then 
% by \textsc{S-Once} we have  
%  \tcHC{\Euid}{\Huid}{C$_0$} and \istay{\Euid}{C$_{E_s}$} and
%     \istay{\Huid}{C$_{H_s}$} and
%     \combine{\Euid}{\Huid}{\C\,$\wedge$\,C$_0$}{C$'$} and
%     \solveH{\Euid$'$}{\Huid$'$}{$($C$'$ $\wedge$ C$_{E_s}$ $\wedge$ C$_{H_s})$}.
% We proceed by structural induction on \texttt{L}:
% \begin{itemize}
% \item \texttt{L} is a variable \texttt{x}.  Then
%  by \textsc{E-Var}
% \Euid\ maps \texttt{x} to \texttt{v}.  First suppose that 
% that \texttt{v} = \texttt{c} for some constant \texttt{c}.  By 
% \textsc{I-StayOne} and \textsc{I-StayConst} we have that
% \texttt{C$_{E_s}$} contains a weak constraint \texttt{
%   x=c}.  Therefore, \Euid$'$ must contain a value for \texttt{x}, which
% we will call \texttt{v$'$}.  If that value is a constant, then the result
% follows.  Otherwise
% \texttt{v$'$} is either an object value or a reference.  But in that case
% there must be a required constraint in \texttt{C$_0$} that forces the
% solver to choose a new value for \texttt{x} and furthermore one that is
% allowed to be an object value or reference.  We argue that this is not possible.

% By \textsc{T-Constant} we have \tcH{\Euid}{\Huid}{x}{PrimitiveType},
% so by the other typing rules \texttt{x}
% can appear as a direct subexpression in arithmetic, comparison, and combination
% expressions in \texttt{C$_0$}.  If \texttt{x} appears in an arithmetic
% expression then \texttt{c} must be a numeric value and any solution will
% also make \texttt{x} numeric.  If \texttt{x} appears in a comparison
% expression then the associated comparison operator was appropriate for
% \texttt{c} and any solution for \texttt{x} satisfying the associated
% constraint must also assign \texttt{x} to a primitive value.  If \texttt{x}
% appears in a combination constraint then \texttt{x} is a boolean value
% and any solution will also make \texttt{x} a boolean.  Notably, \texttt{x}
% cannot appear in an expression of the form \texttt{x.l}, because that
% will not typecheck according to \textsc{T-Field}.  Finally, \texttt{x}
% can appear as a direct subexpression in an object value expression.
% But by \texttt{T-Compare} and \texttt{T-Field} this record and its fields
% can also only be constrained in a manner that preserves \texttt{
%   PrimitiveType}.

% Second, suppose that \texttt{v} is an object \texttt{
%   \{l$_1$:v$_1$,$\ldots$,l$_n$:v$_n$\}}.  Then
% by \textsc{IStayOne} and \textsc{IStayObject} we have that
% \texttt{C$_{E_s}$} contains a required constraint
%  \texttt{x=\{l$_1$:x$_1$,$\ldots$,l$_n$:x$_{n}$\}} and the result
%  follows.  Finally, suppose \texttt{v} is a
% reference \texttt{r}  such that \Deref(\texttt{r}) is an object literal
%  \texttt{\{l$_1$:v$_1$,$\ldots$,l$_n$:v$_n$\}}.  By
% \textsc{IStayRef} we have that 
% \texttt{C$_{E_s}$} contains a required constraint \texttt{x = r}.  
% Further, by \textsc{IStayHeap} we have that
% \texttt{C$_{H_s}$} contains a required constraint 
% of the form  \texttt{x=\{l$_1$:x$_1$,$\ldots$,l$_n$:x$_{n}$\}}.
% Therefore the result follows.

% \item \texttt{L} has the form \texttt{L$'$.l}.  Then the lemma holds
%   inductively for \texttt{L$'$}.  Since \evalH{\Euid}{\Huid}{L}{v}, by
%   the evaluation rules we have two cases.
% First suppose that \textsc{E-ValueField} is the last rule in the
% derivation of \evalH{\Euid}{\Huid}{L}{v}.  Then   
%  \evalH{\Euid}{\Huid}{L$'$}{\{l$_{01}$:v$_{01}$,$\ldots$,l$_{0p}$:v$_{0p}$\}}
%  and \texttt{l} has the form \texttt{l$_i$} and \texttt{v} = \texttt{v$_{0i}$} and 
%   1 $\leq$ $i$ $\leq$ $p$.  By induction   
%  \evalH{\Euid$'$}{\Huid$'$}{L$'$}{\{l$_{01}$:v$_{01}'$,$\ldots$,l$_{0p}$:v$_{0p}'$\}}
%  so by \textsc{E-ValueField} we have
%  \evalH{\Euid}{\Huid}{L}{v$_{0i}'$}.  Similarly, suppose that
%  \textsc{E-Field} is the last rule in the  
% derivation of \evalH{\Euid}{\Huid}{L}{v}.  Then   
%  \evalH{\Euid}{\Huid}{L$'$}{r} and \texttt{\Huid}(\texttt{r}) = 
% \texttt{\{l$_{01}$:v$_{01}$,$\ldots$,l$_{0p}$:v$_{0p}$\}}
%  and \texttt{l} has the form \texttt{l$_i$} and
%   1 $\leq$ $i$ $\leq$ $p$ and \texttt{v} = \texttt{v$_{0i}$}.  By induction   
%  \evalH{\Euid$'$}{\Huid$'$}{L$'$}{r} and \texttt{\Huid}(\texttt{r}) = 
% \texttt{\{l$_{01}$:v$_{01}'$,$\ldots$,l$_{0p}$:v$_{0p}'$\}}
%  so by \textsc{E-Field} we have
%  \evalH{\Euid}{\Huid}{L}{v$_{0i}'$}.  So in both cases we simply need
%  to show that \texttt{v$_{0i}$} and \texttt{v$_{0i}'$} are the same type of
%  values and with the same top-level structure.  This result follows by
%  case analysis of the form of \texttt{v$_{0i}$} using an analogous
%  argument to that in the first bullet above.
% \end{itemize}
% \end{proof}

The second theorem formalizes the idea that all solutions to an
assignment will produce structurally equivalent environments and heaps:

\begin{Theorem} (Structural Determinism)
  If
  \emph{\stepH{\Euid}{\Huid}{\C}{\I}{L := e}{\Euid$_1$}{\Huid$_1$}{\C$_1$}{\I$_1$}} and
  \emph{\stepH{\Euid}{\Huid}{\C}{\I}{L := e}{\Euid$_2$}{\Huid$_2$}{\C$_2$}{\I$_2$}} 
  and \emph{\tcH{\Euid}{\Huid}{x}{T$_0$}}, then there exists a type {\tt T}
  such that
\emph{\tcH{\Euid$_1$}{\Huid$_1$}{x}{T}} and
\emph{\tcH{\Euid$_2$}{\Huid$_2$}{x}{T}}.
\end{Theorem}
\begin{proof}
By {\sc S-Asgn} we have     \evalH{\Euid}{\Huid}{e}{v} 
and \stayPrefix{\Euid}{\Huid}{L}{C$_L$} and
\stayPrefix{\Euid}{\Huid}{I}{C$_I$} and
\phase{\Euid}{\Huid}{C$_L$\,$\wedge$\,C$_I$\,$\wedge$\,L=v\,$\wedge$\,I}{weak}{\Euid$_1'$;\Huid$_1'$} and
    \tcHC{\Euid$_1'$}{\Huid$_1'$}{C} and
\phase{\Euid$_1'$}{\Huid$_1'$}{\C\,$\wedge$\,L=v}{required}{\Euid$_1$;\Huid$_1$}.
Also by {\sc S-Asgn} and Lemma~\ref{lem:det}
we have
\phase{\Euid}{\Huid}{C$_L$\,$\wedge$\,C$_I$\,$\wedge$\,L=v\,$\wedge$\,I}{weak}{\Euid$_2'$;\Huid$_2'$} and
    \tcHC{\Euid$_2'$}{\Huid$_2'$}{C} and
\phase{\Euid$_2'$}{\Huid$_2'$}{\C\,$\wedge$\,L=v}{required}{\Euid$_2$;\Huid$_2$}.
By Lemma~\ref{lem:phase1Det} we have that  {\tt \Euid$_1'$} = {\tt \Euid$_2'$} and {\tt \Huid$_1'$} = {\tt
  \Huid$_2'$}.  Since \tcH{\Euid}{\Huid}{x}{T$_0$}, by {\sc
  T-Variable} {\tt x} is in the domain of {\tt \Euid}, so it is also in
the domain of {\tt \Euid$_1'$} and by {\sc E-Var} we have
\evalH{\Euid$_1'$}{\Huid$_1'$}{x}{v$_x$} where {\tt \Euid$_1'$}({\tt
  x}) = {\tt v$_x$}.
Then by Lemma~\ref{lem:valueImpliesType} there is some type {\tt T$'$}
such that  \tcH{\Euid$_1'$}{\Huid$_1'$}{v$_x$}{T$'$}
and \tcH{\Euid$_1'$}{\Huid$_1'$}{x}{T$'$}.  Finally by
Lemma~\ref{lem:keyAux} we have that \tcH{\Euid$_1$}{\Huid$_1$}{x}{T$'$} and
\tcH{\Euid$_2$}{\Huid$_2$}{x}{T$'$}, so the result is shown with {\tt T} =
{\tt T$'$}.

% As described earlier the rule \textsc{S-Asgn} encodes the two-phase
% process for solving constraints.  We are given two derivations of an
% assignment's execution using this rule:
%   \stepH{\Euid}{\Huid}{\C}{\I}{L := e}{\Euid$_1$}{\Huid$_1$}{\C$_1$}{\I$_1$} and
%   \stepH{\Euid}{\Huid}{\C}{\I}{L :=
%     e}{\Euid$_2$}{\Huid$_2$}{\C$_2$}{\I$_2$}.
% By the lemma below, these two derivations have the identical
% environment and heap after the first phase, which solves identity
% constraints.  Let's call this environment and heap \Euid$'$ and
% \Huid$'$.  Further, the second phase, which solves value constraints, 
% is identical to how a \texttt{
%   once} constraint is processed by the rule \textsc{S-Once}.
% Therefore by the same argument as in the Structure Preservation
% Theorem, if \tcH{\Euid$'$}{\Huid$'$}{x}{T}, then also
% \tcH{\Euid$_1$}{\Huid$_1$}{x}{T} and
% \tcH{\Euid$_2$}{\Huid$_2$}{x}{T}.  By construction \Euid$'$ includes
% all the variables in \Euid, so the result follows.
\end{proof}


\begin{Definition}
\label{def:alias}
We say that {\tt L$_1$} and {\tt L$_2$} are \emph{aliases} given \Euid\
and \Huid\ if either:
\begin{itemize}
\item {\tt L$_1$} and {\tt L$_2$} are the same variable {\tt x}
\item {\tt L$_1$} has the form {\tt L$_1'$.l} and {\tt L$_2$} has the
  form {\tt L$_2'$.l} and there is a reference {\tt r} such that \emph{\evalH{\Euid}{\Huid}{L$_1'$}{r}} and \emph{\evalH{\Euid}{\Huid}{L$_2'$}{r}}
\end{itemize}
\end{Definition}

\begin{Definition}
\label{def:ops}
We say that {\tt L} and {\tt L$'$} are the \emph{operands} of the
constraint {\tt L == L$'$}.
\end{Definition}

\begin{Definition}
\label{def:graph}
We define the \emph{induced graph} of {\tt I} and {\tt L} given \Euid\
and \Huid\ as follows.  Let $S$ be the set that includes {\tt L} as
well as all operands of identity tests in {\tt I}.  Partition $S$ into
equivalence classes defined by the alias relation:  {\tt L$_1$} and {\tt
  L$_2$} are in the same equivalence class if and only if they are
aliases given \Euid\ and \Huid.  % (It is easy to show that the alias
% relation is an equivalence relation.)
Then the induced graph has one node per equivalence class and an
undirected edge between nodes $N_1$ and $N_2$ if there is a conjunct {\tt
  L$_1$==L$_2$} in {\tt I} such that {\tt L$_1$} belongs to node $N_1$
and {\tt L$_2$} belongs to node $N_2$.
\end{Definition}

\begin{Definition}
\label{def:relevant}
We say that a node in the induced graph of {\tt I} and {\tt L} given \Euid\
and \Huid\ is \emph{relevant} if it
is reachable from {\tt L}'s node in the graph; an l-value is
\emph{relevant} if 
its node in the graph is relevant.  % An expression of the form {\tt
  % r.l} is \emph{relevant} if {\tt L.l} is relevant and \emph{\evalH{\Euid}{\Huid}{L}{r}}.
\end{Definition}

\begin{Definition}
\label{def:update}
The \emph{relevant update} of \Euid\ and \Huid\ for {\tt I} and {\tt
  L=v} is the environment \Euid$'$ and heap \Huid$'$ that are identical to \Euid\ and
\Huid\ except that for each relevant l-value {\tt L$_0$} in the induced graph of {\tt I} and {\tt L} given \Euid\
and \Huid:
\begin{itemize}
\item If {\tt L$_0$} is a variable {\tt x}, then \Euid$'$({\tt x}) =
  {\tt v}.
\item If {\tt L$_0$} has the form {\tt L$_0'$.l} and
  \evalH{\Euid}{\Huid}{{\tt L$_0'$}}{{\tt r}},
then \evalH{\Euid$'$}{\Huid$'$}{{\tt r.l}}{{\tt v}}.
\end{itemize}
\end{Definition}

\begin{Lemma}
\label{lem:phase1Det}
If \emph{\stayPrefix{\Euid}{\Huid}{L}{C$_L$}} and
\emph{\stayPrefix{\Euid}{\Huid}{I}{C$_I$}} and 
\emph{\phase{\Euid}{\Huid}{C$_L$\,$\wedge$\,C$_I$\,$\wedge$\,L=v\,$\wedge$\,I}{weak}{\Euid$_1$;\Huid$_1$}}
and
\emph{\phase{\Euid}{\Huid}{C$_L$\,$\wedge$\,C$_I$\,$\wedge$\,L=v\,$\wedge$\,I}{weak}{\Euid$_2$;\Huid$_2$}},
then {\tt \Euid$_1$} = {\tt \Euid$_2$} and {\tt \Huid$_1$} = {\tt
  \Huid$_2$}.
\end{Lemma}
\begin{proof}
By {\sc Solve} we have \gstay{\Euid}{{\tt weak}}{C$_{E}$} and
    \gstay{\Huid}{{\tt weak}}{C$_{H}$} and
    \combine{\Euid}{\Huid}{$($C$_L$\,$\wedge$\,C$_I$\,$\wedge$\,L=v\,$\wedge$\,I$)$}{C$'$}
and \solveH{\Euid$'$}{\Huid$'$}{$($C$'$ $\wedge$ C$_{E}$ $\wedge$
  C$_{H}$$)$} and
\solveH{\Euid$''$}{\Huid$''$}{$($C$'$ $\wedge$ C$_{E}$ $\wedge$
  C$_{H}$$)$}.  

By
Lemma~\ref{lem:key1}, {\tt \Euid$'$} and {\tt \Huid$'$} include all
the updates of the relevant update (which we will call {\tt \Euid$_0$}
and {\tt \Huid$_0$}) of \Euid\ and \Huid\ for {\tt I} and {\tt L=v}, and similarly for {\tt \Euid$''$}
and {\tt \Huid$''$}.  To complete the proof we argue that both of
these solutions are in fact identical to the relevant update.  WLOG we
consider {\tt \Euid$'$} and {\tt \Huid$'$}.  By Lemma~\ref{lem:relevantSolution} the relevant
update of \Euid\ and \Huid\ for {\tt I} and {\tt L=v} is a solution to
the constraint {\tt C$'$ $\wedge$ C$_{E}$ $\wedge$ C$_{H}$}.
Therefore if {\tt \Euid$'$} and {\tt \Huid$'$} is not the relevant
update, then by Definition~\ref{def:update} either:
\begin{itemize}
\item there is a variable {\tt x} such that {\tt \Euid$_0$}({\tt x}) =
  {\tt \Euid}({\tt x}) but {\tt \Euid$'$}({\tt x}) has a different
  value
\item there is a reference {\tt r} and field label {\tt l} such that
  {\tt \Huid$_0$}({\tt r}){\tt .l} =
  {\tt \Huid}({\tt r}){\tt .l} but {\tt \Huid$'$}({\tt r}){\tt .l} has a different
  value
\end{itemize}
Consider the former.  
Since  \gstay{\Euid}{{\tt weak}}{C$_{E}$},
    by {\sc StayEnv}, {\sc StayConst}, and {\sc StayRef} we have that
    {\tt C$_E$} includes a weak constraint {\tt x=v$_x$}, where {\tt
      \Euid}({\tt x}) = {\tt v$_x$}.  Since {\tt \Euid$'$} and {\tt \Huid$'$} include all
the updates of the relevant update, {\tt \Euid$'$} and {\tt \Huid$'$}
satisfy strictly fewer weak constraints than the relevant update,
contradicting the optimality of {\tt \Euid$'$} and {\tt \Huid$'$}.

Similarly, consider the latter.  Since  \gstay{\Huid}{{\tt weak}}{C$_{H}$},
    by {\sc StayHeap} and {\sc StayObject}
    {\tt C$_H$} includes a required constraint {\tt \Deref(r)=\{l$_1$:x$_1$,$\ldots$,l$_n$:x$_n$\}} where the {\tt
    x$_i$} variables are fresh and {\tt l} is some {\tt l$_i$}.   Then
  by  {\sc StayConst} and {\sc StayRef} there is a weak constraint of
  the form {\tt x$_i$ = v$_i$}
where {\tt
      \Huid}({\tt r}){\tt .l$_i$} = {\tt v$_i$}.  Since {\tt \Euid$'$} and {\tt \Huid$'$} include all
the updates of the relevant update, {\tt \Euid$'$} and {\tt \Huid$'$}
satisfy strictly fewer weak constraints than the relevant update,
contradicting the optimality of {\tt \Euid$'$} and {\tt \Huid$'$}.

\end{proof}

% In the following lemma, we implicitly convert a constraint {\tt C} to an
% expression by removing all priorities, thereby allowing the constraint
% to be evaluated.

\begin{Lemma}
\label{lem:relevantSolution}
If    \emph{\gstay{\Euid}{{\tt weak}}{C$_{E}$}} and
    \emph{\gstay{\Huid}{{\tt weak}}{C$_{H}$}} and
\emph{\stayPrefix{\Euid}{\Huid}{L}{C$_L$}} and
\emph{\stayPrefix{\Euid}{\Huid}{I}{C$_I$}} and
    \emph{\combine{\Euid}{\Huid}{$($C$_L$\,$\wedge$\,C$_I$\,$\wedge$\,L=v\,$\wedge$\,I$)$}{C$'$}}
    and the constraint {\tt C$'$ $\wedge$ C$_{E}$ $\wedge$ C$_{H}$} is satisfiable,
then the 
relevant update {\tt \Euid$'$} and {\tt \Huid$'$} of \Euid\ and \Huid\ for {\tt I} and {\tt L=v} 
is a solution to the constraint {\tt C$'$ $\wedge$ C$_{E}$ $\wedge$ C$_{H}$}.
\end{Lemma}
\begin{proof}
It suffices to show that all required constraints in {\tt C$'$
  $\wedge$ C$_{E}$ $\wedge$ C$_{H}$} are satisfied in  {\tt \Euid$'$}
and {\tt \Huid$'$}.  We consider the various constraints in
turn:
\begin{itemize}
\item {\tt C$_E$}:  Since \gstay{\Euid}{{\tt weak}}{C$_{E}$}, by {\sc
    StayEnv}, {\sc StayConst}, and {\sc StayRef} there are no required
  constraints in {\tt C$_E$}, so all required constraints are
  satisfied vacuously.

\item {\tt C$_H$}:  Since \gstay{\Huid}{{\tt weak}}{C$_{H}$}, by {\sc
    StayHeap}, {\sc StayObject}, {\sc StayConst}, and {\sc StayRef}
  the only required
  constraints in {\tt C$_H$} have the form {\tt
    \Deref(r)=\{l$_1$:x$_1$,$\ldots$,l$_n$:x$_n$\}} where the {\tt
    x$_i$} variables are fresh and {\tt \Huid} maps {\tt r} to some value of the
  form {\tt \{l$_1$:v$_1$,$\ldots$,l$_n$:v$_n$\}}.  By
  Definition~\ref{def:update} also {\tt \Huid$'$} maps {\tt r} to a value of the
  form {\tt \{l$_1$:v$_1'$,$\ldots$,l$_n$:v$_n'$\}} so the constraint
  is satisfied.

\item {\tt C$_L$}:  By {\sc StayPrefixField} the conjuncts in C$_L$
  have the form {\tt x=v} or {\tt r.l=v}.  Suppose the relevant update fails to
  satisfy one of these conjuncts.  We consider each form in turn:
\begin{itemize}
\item {\tt x=v}:  Then the relevant update maps {\tt x} to some {\tt
    v$'$} $\neq$ {\tt v} in {\tt \Euid$'$}.  But by {\sc Solve} and
  Lemma~\ref{lem:key1} any solution to  the constraint {\tt C$'$
    $\wedge$ C$_{E}$ $\wedge$ C$_{H}$} must map {\tt x} to {\tt v$'$}
  in the environment, which violates the constraint {\tt x=v}.  So there must be no solution to the constraints,
  contradicting our assumption of satisfiability.

\item {\tt r.l=v}:  Then the relevant update maps {\tt r.l} to some {\tt
    v$'$} $\neq$ {\tt v} in {\tt \Huid$'$}.  But by {\sc Solve} and
  Lemma~\ref{lem:key1} any solution to  the constraint {\tt C$'$
    $\wedge$ C$_{E}$ $\wedge$ C$_{H}$} must map {\tt r.l} to {\tt v$'$}
  in the environment, which violates the constraint {\tt r.l=v}.  So there must be no solution to the constraints,
  contradicting our assumption of satisfiability.
\end{itemize}

\item {\tt C$_I$}:  By {\sc StayPrefixIdent}  the conjuncts in C$_L$
  have the form {\tt x=v} or {\tt r.l=v}.  So the argument is
  identical to that above for the {\tt C$_L$} constraint.

\item {\tt L=v}:  By Definitions~\ref{def:graph} and \ref{def:update}
  we have that {\tt L} is relevant for {\tt I} and {\tt L} given
  \Euid\ and \Huid.  Then by Lemma~\ref{lem:key2} we have 
\evalH{\Euid$'$}{\Huid$'$}{L}{v}, so by {\sc E-Compare} we have
\evalH{\Euid$'$}{\Huid$'$}{L=v}{true}.

\item {\tt I}:  Consider a conjunct {\tt L$_0$ == L$_1$} in {\tt I}.
  We have two cases:
\begin{itemize}
\item {\tt L$_0$} is relevant for {\tt I} and {\tt L} given \Euid\ and
  \Huid.  By Definitions~\ref{def:relevant} and \ref{def:graph}, {\tt
    L$_1$} is also relevant.  Then by Lemma~\ref{lem:key2} we have
\evalH{\Euid$'$}{\Huid$'$}{L$_0$}{v} and
\evalH{\Euid$'$}{\Huid$'$}{L$_1$}{v}, so by {\sc E-IdentityTrue} we
have
\evalH{\Euid$'$}{\Huid$'$}{L$_0$==L$_1$}{true}.

\item {\tt L$_0$} is not relevant for {\tt I} and {\tt L} given \Euid\ and
  \Huid.  By Definitions~\ref{def:relevant} and \ref{def:graph}, {\tt
    L$_1$} is also not relevant.  By well formedness we have
\evalH{\Euid}{\Huid}{L$_0$==L$_1$}{true}, so by {\sc E-IdentityTrue}
there is some {\tt v$_0$} such that
\evalH{\Euid}{\Huid}{L$_0$}{v$_0$} and
\evalH{\Euid}{\Huid}{L$_1$}{v$_0$}.
Then by Lemma~\ref{lem:key3} we have
\evalH{\Euid$'$}{\Huid$'$}{L$_0$}{v$_0$} and
\evalH{\Euid$'$}{\Huid$'$}{L$_1$}{v$_0$}, so by {\sc E-IdentityTrue} we
have
\evalH{\Euid$'$}{\Huid$'$}{L$_0$==L$_1$}{true}. 
\end{itemize}

\end{itemize}
\end{proof}

\begin{Lemma}
\label{lem:key1}
If \emph{\stayPrefix{\Euid}{\Huid}{L}{C$_L$}} and
\emph{\stayPrefix{\Euid}{\Huid}{I}{C$_I$}} and
\emph{\phase{\Euid}{\Huid}{C$_L$\,$\wedge$\,C$_I$\,$\wedge$\,L=v\,$\wedge$\,I}{weak}{\Euid$'$;\Huid$'$}}
and {\tt L$_0$} is a relevant l-value of {\tt I} and {\tt L} given
\Euid\ and \Huid, then \emph{\evalH{\Euid$'$}{\Huid$'$}{{\tt
      L$_0$}}{{\tt v}}} and:
\begin{itemize}
\item If {\tt L$_0$} is a variable {\tt x}, then \Euid$'$({\tt x}) =
  {\tt v}.
\item If {\tt L$_0$} has the form {\tt L$_0'$.l} and
  \evalH{\Euid}{\Huid}{{\tt L$_0'$}}{{\tt r}},
then \evalH{\Euid$'$}{\Huid$'$}{{\tt r.l}}{{\tt v}}.
\end{itemize}
\end{Lemma}
\begin{proof}

By Definition~\ref{def:relevant}, we know that {\tt L$_0$}'s node in the induced
graph for {\tt I} and {\tt L} given \Euid\ and \Huid\ is reachable
from {\tt L}'s node.  The proof proceeds by
induction on the length $k$ of the path between these nodes.
\begin{itemize}
\item Case $k=0$.  Then by Definition~\ref{def:graph}, {\tt L} and {\tt L$_0$}
  are aliases given \Euid\ and \Huid.  Case analysis on the structure
  of {\tt L}:
  \begin{itemize}
  \item Case {\tt L} is a variable {\tt x}.  Then by
  Definition~\ref{def:alias} also
 {\tt L$_0$} is {\tt x}.  Since
 \phase{\Euid}{\Huid}{C$_L$\,$\wedge$\,C$_I$\,$\wedge$\,L=v\,$\wedge$\,I}{weak}{\Euid$'$;\Huid$'$}
 we must have \evalH{\Euid$'$}{\Huid$'$}{x=v}{true}, so by {\sc
   E-Compare} and the semantics of equality we have
 \evalH{\Euid$'$}{\Huid$'$}{x}{v}.  Then by
 {\sc E-Var} also \Euid$'$({\tt x}) =
  {\tt v}.

  \item Case {\tt L} has the form {\tt L$'$.l}.  Then by
    Definition~\ref{def:alias} also {\tt L$_0$} has the
    form {\tt L$_0'$.l} and \evalH{\Euid}{\Huid}{L$'$}{r}
    and \evalH{\Euid}{\Huid}{L$_0'$}{r}.  
We are given 
\stayPrefix{\Euid}{\Huid}{L}{C$_L$}. Then since \phase{\Euid}{\Huid}{C$_L$\,$\wedge$\,C$_I$\,$\wedge$\,L=v\,$\wedge$\,I}{weak}{\Euid$'$;\Huid$'$}
 we must have \evalH{\Euid$'$}{\Huid$'$}{C$_{L}$}{true}.
Similarly we are given 
\stayPrefix{\Euid}{\Huid}{I}{C$_I$} so by {\sc StayPrefixIdent} we
have also \stayPrefix{\Euid}{\Huid}{L$_0$}{C$_{L_0}$} where {\tt
  C$_{L_0}$} is a conjunct within {\tt C$_I$}.  Then since \phase{\Euid}{\Huid}{C$_L$\,$\wedge$\,C$_I$\,$\wedge$\,L=v\,$\wedge$\,I}{weak}{\Euid$'$;\Huid$'$}
 we must have \evalH{\Euid$'$}{\Huid$'$}{C$_{L_0}$}{true}.  So by
Lemma~\ref{lem:det} and Lemma~\ref{lem:stayPrefix} we have 
\evalH{\Euid$'$}{\Huid$'$}{L$'$}{r} and
\evalH{\Euid$'$}{\Huid$'$}{L$_0'$}{r}.  
 We also know \evalH{\Euid$'$}{\Huid$'$}{L=v}{true}, so by {\sc
   E-Compare} and the semantics of equality we have
 \evalH{\Euid$'$}{\Huid$'$}{L}{v}.  Then by
 {\sc E-Field} we have \Huid$'$({\tt r}) is a record whose {\tt l}
 field has value {\tt v}.  Then by {\sc E-Field} also \evalH{\Euid$'$}{\Huid$'$}{{\tt r.l}}{{\tt v}}
and \evalH{\Euid$'$}{\Huid$'$}{{\tt L$_0$}}{{\tt v}}.

  \end{itemize}

\item Case $k>0$.  Then by Definition~\ref{def:graph} there is some
  neighbor node of {\tt L$_0$}'s node that is only $k-1$ hops away
  from {\tt L}'s node, which contains an l-value {\tt L$_1'$} such
  that {\tt L$_0'$ == L$_1'$} or vice versa is in {\tt I}, where {\tt
    L$_0'$} is an alias of {\tt L$_0$} given \Euid\ and \Huid.  
By induction we have \evalH{\Euid$'$}{\Huid$'$}{L$_1'$}{v}.  Then
since the identities in {\tt I} are satisfied in {\tt \Euid$'$} and
{\tt \Huid$'$}, by {\sc E-IdentityTrue} also
\evalH{\Euid$'$}{\Huid$'$}{L$_0'$}{v}.  Case
analysis on the structure of {\tt L$_0'$}:
  \begin{itemize}

 \item Case {\tt L$_0'$} is a variable {\tt x}.  Then by
  Definition~\ref{def:alias} also
 {\tt L$_0$} is {\tt x} and by  {\sc E-Var} also \Euid$'$({\tt x}) =
  {\tt v}.

  \item Case {\tt L$_0'$} has the form {\tt L$'$.l}.  Then by
    Definition~\ref{def:alias} also {\tt L$_0$} has the
    form {\tt L$_{00}$.l} and \evalH{\Euid}{\Huid}{L$'$}{r}
    and \evalH{\Euid}{\Huid}{L$_{00}$}{r}.  
We are given 
\stayPrefix{\Euid}{\Huid}{I}{C$_I$} so by {\sc StayPrefixIdent} we
have also \stayPrefix{\Euid}{\Huid}{L$_0'$}{C$_{L_0'}$} where {\tt
  C$_{L_0'}$} is a conjunct within {\tt C$_I$}.  Then since \phase{\Euid}{\Huid}{C$_L$\,$\wedge$\,C$_I$\,$\wedge$\,L=v\,$\wedge$\,I}{weak}{\Euid$'$;\Huid$'$}
 we must have \evalH{\Euid$'$}{\Huid$'$}{C$_{L_0'}$}{true}.  By the
 same argument we must also have
 \evalH{\Euid$'$}{\Huid$'$}{C$_{L_0}$}{true} where \stayPrefix{\Euid}{\Huid}{L$_0$}{C$_{L_0}$}.  So by
 Lemma~\ref{lem:det} and Lemma~\ref{lem:stayPrefix} we have 
\evalH{\Euid$'$}{\Huid$'$}{L$'$}{r} and
\evalH{\Euid$'$}{\Huid$'$}{L$_{00}$}{r}.  
Then since \evalH{\Euid$'$}{\Huid$'$}{L$_0'$}{v} by
 {\sc E-Field} we have \Huid$'$({\tt r}) is a record whose {\tt l}
 field has value {\tt v}.  Then by {\sc E-Field} also \evalH{\Euid$'$}{\Huid$'$}{{\tt r.l}}{{\tt v}}
and \evalH{\Euid$'$}{\Huid$'$}{{\tt L$_0$}}{{\tt v}}.


\end{itemize}

\end{itemize}
\end{proof}


\begin{Lemma}
\label{lem:key2}
If \emph{\stayPrefix{\Euid}{\Huid}{L}{C$_L$}} and
\emph{\stayPrefix{\Euid}{\Huid}{I}{C$_I$}} and
{\tt \Euid$'$} and {\tt \Huid$'$} is the relevant update of \Euid\
and \Huid\ for {\tt I} and {\tt L=v} 
and \emph{\evalH{\Euid$'$}{\Huid$'$}{C$_L$}{true}} and
and \emph{\evalH{\Euid$'$}{\Huid$'$}{C$_I$}{true}} and {\tt L$_0$} is
a relevant l-value for {\tt I} and {\tt L} given \Euid\ and \Huid, then
\emph{\evalH{\Euid$'$}{\Huid$'$}{L$_0$}{v}}.
\end{Lemma}
\begin{proof}
Case analysis of the structure of {\tt L$_0$}:
\begin{itemize}
\item Case {\tt L$_0$} is a variable {\tt x}:  By Definition~\ref{def:update}
{\tt \Euid$'$}({\tt x}) = {\tt v}, so the result follows by {\sc
  E-Var}.
\item Case {\tt L$_0$} has the form {\tt L$'$.l}:  First we argue that
  \stayPrefix{\Euid}{\Huid}{L$_0$}{C$_{L0}$} and
  \evalH{\Euid$'$}{\Huid$'$}{C$_{L0}$}{true}.  If {\tt L$_0$} is {\tt
    L}, then these follow from the assumptions of the lemma.
  Otherwise, by Definition~\ref{def:relevant} {\tt L$_0$} is an
  operand in an identity test in {\tt I}.  Then since
  \stayPrefix{\Euid}{\Huid}{I}{C$_I$} by  {\sc StayPrefixIdent} we
  have \stayPrefix{\Euid}{\Huid}{L$_0$}{C$_{L0}$}, where {\tt
    C$_{L0}$} is a conjunct in {\tt C$_I$}.  Then since
  \evalH{\Euid$'$}{\Huid$'$}{C$_I$}{true} also
  \evalH{\Euid$'$}{\Huid$'$}{C$_{L0}$}{true}.

Then by
  Lemma~\ref{lem:stayPrefix}
 there is some {\tt r$'$} such that
\evalH{\Euid}{\Huid}{L$'$}{r$'$} and
\evalH{\Euid$'$}{\Huid$'$}{L$'$}{r$'$}.  By Lemma~\ref{lem:det}
and Definition~\ref{def:update} we have \evalH{\Euid}{\Huid}{{\tt
    r$'$.l}}{{\tt v}} and the result follows by {\sc E-Field}.
\end{itemize}
\end{proof}

\begin{Lemma}
\label{lem:key3}
If \emph{\stayPrefix{\Euid}{\Huid}{L}{C$_L$}} and
\emph{\stayPrefix{\Euid}{\Huid}{I}{C$_I$}} and
{\tt \Euid$'$} and {\tt \Huid$'$} is the relevant update of \Euid\
and \Huid\ for {\tt I} and {\tt L=v} 
and \emph{\evalH{\Euid$'$}{\Huid$'$}{C$_L$}{true}} and
and \emph{\evalH{\Euid$'$}{\Huid$'$}{C$_I$}{true}} and {\tt L$_0$} is
an operand of an identity test in {\tt I} and {\tt L$_0$} is not
relevant for {\tt I} and {\tt L} given \Euid\ and \Huid\ and \emph{\evalH{\Euid}{\Huid}{L$_0$}{v$_0$}},
\emph{\evalH{\Euid$'$}{\Huid$'$}{L$_0$}{v$_0$}}.
\end{Lemma}
\begin{proof}
Case analysis on the structure of {\tt L$_0$}:
\begin{itemize}
\item Case {\tt L$_0$} is a variable {\tt x}.  Since {\tt x} is not
  relevant, by
  Definition~\ref{def:update} the value of {\tt x} in {\tt \Euid$'$}
  is the same as that in {\tt Euid}.  Since
  \evalH{\Euid}{\Huid}{L$_0$}{v$_0$}, by {\sc E-Var} we have
  \Euid({\tt x}) = {\tt v$_0$}, so also {\tt \Euid$'$}({\tt x}) = {\tt
    v$_0$} and the result follows by {\sc E-Var}.

\item Case {\tt L$_0$} has the form {\tt L$'$.l}.   Since {\tt L$_0$}
  is an operand in an identity test in {\tt I} and
  \stayPrefix{\Euid}{\Huid}{I}{C$_I$}, by  {\sc StayPrefixIdent} we
  have \stayPrefix{\Euid}{\Huid}{L$_0$}{C$_{L0}$}, where {\tt
    C$_{L0}$} is a conjunct in {\tt C$_I$}.  Then since
  \evalH{\Euid$'$}{\Huid$'$}{C$_I$}{true} also
  \evalH{\Euid$'$}{\Huid$'$}{C$_{L0}$}{true}.

Therefore by
  Lemma~\ref{lem:stayPrefix}
 there is some {\tt r$'$} such that
\evalH{\Euid}{\Huid}{L$'$}{r$'$} and
\evalH{\Euid$'$}{\Huid$'$}{L$'$}{r$'$}.  Since
\evalH{\Euid}{\Huid}{L$_0$}{v$_0$} by {\sc E-Field} and Lemma~\ref{lem:det} we have that \Deref({\tt r$'$}){\tt .l} is {\tt v$_0$}.  By
Definition~\ref{def:update} {\tt \Huid$'$}({\tt r$'$}) also has a
field with label {\tt l}, so by {\sc E-Field} we are done if that
field's value is {\tt v$_0$}.  Suppose not.  Then by
Definition~\ref{def:update} there is some relevant l-value {\tt L$_1$}
of the form {\tt L$_1'$.l} such that
\evalH{\Euid}{\Huid}{L$_1'$}{r$'$}.  But then by
Definition~\ref{def:alias} we have that {\tt L$_0$} and {\tt L$_1$}
are aliases so they correspond to the same node in the induced graph of {\tt I} and {\tt L} given \Euid\
and \Huid\
by Definition~\ref{def:graph}.  But then since {\tt L$_1$} is
relevant, by Definition~\ref{def:relevant} so is {\tt L$_0$} and we
have a contradiction.
\end{itemize}
\end{proof}


\begin{Lemma}
\label{lem:stayPrefix}
If \emph{\stayPrefix{\Euid}{\Huid}{L.f}{C}} and
\emph{\evalH{\Euid$'$}{\Huid$'$}{C}{true}}, then there is some
reference {\tt r} such that 
\emph{\evalH{\Euid}{\Huid}{L}{r}} and
\emph{\evalH{\Euid$'$}{\Huid$'$}{L}{r}}.
\end{Lemma}
\begin{proof}
By {\sc StayPrefixField} we have {\tt L.f} = {\tt
x.l$_1$.$\ldots$.l$_n$} and $n>0$ and \evalH{\Euid}{\Huid}{x}{r}
and \evalH{\Euid}{\Huid}{x.l$_1$}{r$_1$} $\cdots$
\evalH{\Euid}{\Huid}{x.l$_1$.$\ldots$.l$_{n-1}$}{r$_{n-1}$} and {\tt
  C} is {\tt
  x=r\,$\wedge$\,x.l$_1$=r$_1$\,$\wedge$\,$\cdots$\,$\wedge$\,x.l$_1$.$\ldots$.l$_{n-1}$=r$_{n-1}$}.  
\begin{itemize}
\item Case $n=1$.  Then {\tt L} is {\tt x} and {\tt C} is {\tt x=r}.
  Since \evalH{\Euid$'$}{\Huid$'$}{C}{true}, by {\sc E-Compare} and
  the semantics of equality we have 
  \evalH{\Euid$'$}{\Huid$'$}{x}{v$_1$} and
  \evalH{\Euid$'$}{\Huid$'$}{r}{v$_2$} and {\tt v$_1$} = {\tt v$_2$}.  By
{\sc E-Ref} we have that {\tt v$_2$} = {\tt r}, so the result follows.
\item Case $n>1$.  Then {\tt L} is {\tt x.l$_1$.$\ldots$.l$_{n-1}$}.
  Since \evalH{\Euid$'$}{\Huid$'$}{C}{true}, by {\sc E-Combine}, {\sc E-Compare}, and
  the semantics of equality we have 
  \evalH{\Euid$'$}{\Huid$'$}{L}{v$_1$} and
  \evalH{\Euid$'$}{\Huid$'$}{r$_{n-1}$}{v$_2$} and {\tt v$_1$} = {\tt v$_2$}.  By
{\sc E-Ref} we have that {\tt v$_2$} = {\tt r$_{n-1}$}, so the result follows.
\end{itemize}
\end{proof}

\begin{Lemma}
\label{lem:det}
(Determinism)
\begin{itemize}
\item If \evalH{\Euid}{\Huid}{e}{v$_1$} and \evalH{\Euid}{\Huid}{e}{v$_2$}
then {\tt v$_1$} = {\tt v$_2$}.
\item If \gstay{\Euid}{$\rho$}{C$_1$} and \gstay{\Euid}{$\rho$}{C$_2$}
  then {\tt C$_1$} = {\tt C$_2$}.
\item If \gstay{\Huid}{$\rho$}{C$_1$} and \gstay{\Huid}{$\rho$}{C$_2$}
  then {\tt C$_1$} = {\tt C$_2$}.
\item If \stayPrefix{\Euid}{\Huid}{L}{C$_1$} and \stayPrefix{\Euid}{\Huid}{L}{C$_2$} 
  then {\tt C$_1$} = {\tt C$_2$}.
\item If \stayPrefix{\Euid}{\Huid}{I}{C$_1$} and \stayPrefix{\Euid}{\Huid}{I}{C$_2$} 
  then {\tt C$_1$} = {\tt C$_2$}.
\end{itemize}
\end{Lemma}
\begin{proof}
Straightforward.
\end{proof}

% \begin{Lemma}
% If  \stay{\Euid}{C$_{E_s}$} and \stay{\Huid}{C$_{H_s}$} and
% \evalH{\Euid}{\Huid}{I}{true} and
%     \combine{\Euid}{\Huid}{\I\,$\wedge$\,L=v}{C$_I$} and
%     \solveH{\Euid$_1$}{\Huid$_1$}{(C$_I$ $\wedge$ C$_{E_s}$ $\wedge$
%       C$_{H_s}$)} and
%     \solveH{\Euid$_2$}{\Huid$_2$}{(C$_I$ $\wedge$ C$_{E_s}$ $\wedge$
%       C$_{H_s}$)}, then \texttt{\Euid$_1$} = \texttt{\Euid$_2$} and
% \texttt{\Huid$_1$} = \texttt{\Huid$_2$}.
% \end{Lemma}
% \begin{proof}
% Let \texttt{C} = \texttt{(C$_I$ $\wedge$ C$_{E_s}$ $\wedge$
%       C$_{H_s}$)}.
% The \textsc{Stay} rules assert weak constraints that the environment
% and heap stay unchanged by the constraint solver.  By assumption
% \evalH{\Euid}{\Huid}{I}{true} as well.  Therefore, in the absence of
% the constraint \texttt{L=v} there is a unique optimal solution to \texttt{C},
% which satisfies all the required constraints and as many 
% weak constraints as possible.  In this case, all the weak constraints
% are satisfied, and the solution is simply to leave
% the environment and heap unchanged.  If
% \evalH{\Euid}{\Huid}{L==v}{true} then \texttt{\Euid} and \texttt{\Huid} are still the
% unique optimal solution to the constraints, so the result follows.

% Otherwise \evalH{\Euid}{\Huid}{L==v}{false}.  The unique optimal
% solution to this constraint in the context of the stay constraints
% (ignoring \texttt{I} for the moment) is the environment and heap \Euid$'$
% and \Huid$'$ that are
% identical to \texttt{\Euid} and \texttt{\Huid} but with \texttt{L} assigned to \texttt{v}.  If
% \evalH{\Euid$'$}{\Huid$'$}{I}{true} then we are done.
% Otherwise, there is at least one constraint in \texttt{I} of the form
% \texttt{L$_1$ = L$_2$} that is violated.  

% We do a subcase analysis, which
% relies on the following definition of the {\em l-value} of
% \texttt{L}.  The {\em l-value} of \texttt{x} under a given \texttt{\Euid} and \texttt{\Huid} 
% is \texttt{x}.  The {\em l-value}
% of \texttt{L.l} under a given \texttt{\Euid} and \texttt{\Huid} is \texttt{v.l}, where
% \evalH{\Euid}{\Huid}{L}{v}.
% Without loss of generality, we
% have three subcases:
% \begin{itemize}
% \item The {\em l-value} of \texttt{L$_1$} under \Euid$'$ and \Huid$'$ 
% is \texttt{L}.  Then there is a unique optimal solution
%   to this constraint in the context of the stay constraints and the
%   constraint \texttt{L=v}, namely the environment and heap \Euid$''$
% and \Huid$''$ that are identical to \Euid$'$ and \Huid$'$ but with \texttt{
%   L$_2$} assigned to \texttt{v}.
% \item \texttt{L$_1$} has the form \texttt{L$_1'$.l$_1$.$\cdots$.l$_n$} where
%   the {\em l-value} of \texttt{L$_1'$} under \Euid$'$ and \Huid$'$ 
% is \texttt{L}.  Then there is a unique optimal solution
%   to this constraint in the context of the stay constraints and the
%   constraint \texttt{L=v}, namely the environment and heap \Euid$''$
% and \Huid$''$ that are identical to \Euid$'$ and \Huid$'$ but with \texttt{
%   L$_2$} assigned to the value of \texttt{v.l$_1$.$\cdots$.l$_n$}.
% \item There is some \texttt{L$_1'$} of the form  \texttt{L$_1$.l$_1$.$\cdots$.l$_n$} such that
%   the {\em l-value} of \texttt{L$_1'$} under \Euid$'$ and \Huid$'$ 
% is \texttt{L}.  Then there is a unique optimal solution
%   to this constraint in the context of the stay constraints and the
%   constraint \texttt{L=v}, namely the environment and heap \Euid$''$
% and \Huid$''$ that are identical to \Euid$'$ and \Huid$'$ but with \texttt{
%   L$_2$.l$_1$.$\cdots$.l$_n$} assigned to the value of \texttt{v}.
% \end{itemize}

% This process continues, identifying constraints in \texttt{I} that are
% falsified, finding the unique optimal solution to such constraints in
% the context of the previously solved constraints and the stay
% constraints.  Since we are given that \texttt{C} is satisfiable,
% eventually this process will terminate with all constraints in \texttt{I}
% satisfied.  Since at each step there was a unique optimal solution, it
% must be the case that \texttt{\Euid$_1$} = \texttt{\Euid$_2$} and
% \texttt{\Huid$_1$} = \texttt{\Huid$_2$}.

% \end{proof}


\end{document}
